\section*{अध्याय \ref{ch:b2:schrodinger-wave-functions} --- श्रोडिंगर समीकरण और तरंग फलन}

\begin{solution}{exo:b2:schrodinger-wave-functions:1}
$h/\sqrt{2mE}$: \qty{0.123}{nm}, \qty{12.3}{pm}, \qty{3.9}{pm} (\qty{100}{keV} पर
गतिज ऊर्जा $mc^2$ की एक-पंचमांश है; आपेक्षिकीय संवेग
\qty{3.7}{pm} देता है); न्यूट्रॉन \qty{0.18}{nm}; C$_{60}$ \qty{2.8}{pm}; टेनिस गेंद
$\sim 10^{-34}\,\unit{m}$। गेंद के सिवा सबने व्यतिकरण दिखाया है।
\end{solution}

\begin{solution}{exo:b2:schrodinger-wave-functions:2}
$A = (\pi a^2)^{-1/4}$; $\langle x\rangle = 0$, $\Delta x = a/\sqrt2$; $0.84$; $0.9995$।
\end{solution}

\begin{solution}{exo:b2:schrodinger-wave-functions:3}
प्रतिस्थापन $E\varphi = -(\hbar^2/2m)\varphi'' + V\varphi$ देता है। $E/h =
\qty{2.4e14}{Hz}$; \qty{4.8e14}{Hz}, \qty{620}{nm}: वही लाल प्रकाश जो संक्रमण
उत्सर्जित करता है।
\end{solution}

\begin{solution}{exo:b2:schrodinger-wave-functions:4}
इलेक्ट्रॉन, \qty{1}{nm}: $\tau = \qty{1.7e-14}{s}$, \qty{1}{ns} के बाद \qty{60}{\micro m},
\qty{1}{s} के बाद \qty{60}{km}; \qty{1}{mm}: $\tau = \qty{17}{s}$, \qty{1}{ns} के बाद अपरिवर्तित,
\qty{1}{s} के बाद $+0.2\%$; प्रोटॉन, \qty{1}{nm}: $\tau = \qty{3.2e-11}{s}$, \qty{30}{nm},
\qty{30}{m}; कण: $\tau = 600$ वर्ष। दृश्य वस्तुओं के लिए $\tau$
खगोलीय है --- और उनका परिवेश उन्हें स्थानीकृत करता रहता है।
\end{solution}

\begin{solution}{exo:b2:schrodinger-wave-functions:5}
(क) \cref{prop:b2:schrodinger-wave-functions:current}। (ख) $|A|^2\hbar k/m$;
$0$; $(|A|^2 - |B|^2)\hbar k/m$। (ग) आपाती और परावर्तित फ्लक्स
बिना किसी व्यतिकरण पद के घट जाते हैं: पारगमन तथा \omterm{thm:b2:wave-interfaces:normal}{परावर्तन
गुणांक} फ्लक्सों से परिभाषित किए जा सकते हैं। (घ) $\psi^*\psi'$ वास्तविक, और उसका
काल्पनिक भाग शून्य।
\end{solution}

\begin{solution}{exo:b2:schrodinger-wave-functions:6}
(क) $A = (2\pi\sigma^2)^{-1/4}$, $\langle x\rangle = 0$, $\Delta x = \sigma$। (ख) $\hbar k_0$।
(ग) $\sigma \times \hbar/2\sigma = \hbar/2$। (घ) $\dd\langle x\rangle/\dd t = \int x\,
\partial_t\rho = -\int x\,\partial_xj = \int j = (\hbar/m)\operatorname{Im}\int\psi^*
\psi' = \langle p\rangle/m$।
\end{solution}

\begin{solution}{exo:b2:schrodinger-wave-functions:7}
(क) \qty{5.9e6}{m/s}, \qty{0.123}{nm}, \qty{51}{ns}। (ख) $v_\varphi = \qty{3.0e6}{m/s}$,
$v_{\text{सम}} = \qty{5.9e6}{m/s}$: बाद वाला। (ग) $\tau = \qty{1.7}{\micro s} \gg
\qty{51}{ns}$: अपरिवर्तित। (घ) $\lambda/a = \qty{4e-7}{rad}$, \qty{0.1}{\micro m}:
किरणें ठीक हैं।
\end{solution}

\begin{solution}{exo:b2:schrodinger-wave-functions:8}
(क) $|\psi|^2$ और हर प्रत्याशा मान अपरिवर्तित रहते हैं। (ख) $\tfrac12(|\psi_1|^2
+ |\psi_2|^2) + \operatorname{Re}(\eu^{\iu\alpha}\psi_1^*\psi_2)$। (ग) दोनों झिरियों से आती
तरंगें; उनकी सापेक्ष कला $k(r_2 - r_1)$ है। (घ) संसूचक
दोनों पथों के लिए भिन्न अवस्थाओं में पहुँचता है; दोनों घटक अब
व्यतिकरण नहीं करते, और प्रतिरूप दो एक-झिरी प्रतिरूपों
का योग है।
\end{solution}

\begin{solution}{exo:b2:schrodinger-wave-functions:9}
(क) $\langle V'\rangle = m\omega^2\langle x\rangle$: ठीक-ठीक। (ख) $\langle V'\rangle
= mg$: ठीक-ठीक। (ग) जब पुंज उस मापक के सामने सँकरा हो जिस पर
$V''$ बदलता है। (घ) \qty{0.1}{m/s} पर उसकी तरंगदैर्घ्य $10^{-29}\,\unit{m}$ है;
मायने रखने के लिए पुंज को कटोरे जितना चौड़ा होना पड़ता। किसी परमाणु की
तापीय तरंगदैर्घ्य, \qty{0.1}{nm}, उसका अपना आकार है।
\end{solution}

\begin{solution}{exo:b2:schrodinger-wave-functions:10}
(क) $\tfrac12(\varphi_1^2 + \varphi_2^2) + \varphi_1\varphi_2\cos\omega_{21}t$ मिलता है। (ख) $L/2 +
\cos\omega_{21}t\int x\varphi_1\varphi_2 = L/2 - (16L/9\pi^2)\cos\omega_{21}t$। (ग)
$3h^2/8mL^2 = \qty{1.13}{eV}$, \qty{2.7e14}{Hz}, \qty{1.1}{\micro m}। (घ) वह
$\nu_{21}$ पर विकिरित करता है: वही रेखा; किसी शुद्ध \omterm{thm:b2:schrodinger-wave-functions:stationary}{स्थायी अवस्था} का कोई दोलायमान
द्विध्रुव नहीं और वह विकिरित नहीं करती।
\end{solution}

\begin{solution}{exo:b2:schrodinger-wave-functions:11}
(क) $\psi \propto \int\eu^{-\sigma_0^2k^2}\eu^{\iu(kx - \hbar k^2t/2m)}\dd k$ लिखिए। (ख) $\alpha =
\sigma_0^2 + \iu\hbar t/2m$, $\beta = \iu x$: $|\psi|^2 \propto \exp[-x^2/2\sigma_0^2(1 +
(\hbar t/2m\sigma_0^2)^2)]$। (ग) $\Delta v = \hbar/2m\sigma_0$, और $\Delta v\,\tau =
\sigma_0$। (घ) $\omega = ck$: कोई प्रकीर्णन नहीं, सारे घटक $c$ पर।
\end{solution}

\begin{solution}{exo:b2:schrodinger-wave-functions:12}
(क) $\Gamma = \hbar/\tau$: \qty{6.6e-8}{eV}; \qty{6.6e-8}{eV}; \qty{66}{MeV}। (ख)
$1/2\pi\tau = \qty{16}{MHz}$, डॉप्लर से सौ गुना कम। (ग) ठंडे
परमाणुओं और पाशित आयनों में डॉप्लर चौड़ाई लुप्त हो जाती है और नैसर्गिक
चौड़ाई तक पहुँचा जाता है --- आवृत्ति मानकों का आधार। (घ) $c\tau =
\qty{3}{m}$, रेल की लंबाई।
\end{solution}

\begin{solution}{pb:b2:schrodinger-wave-functions:1}
\textbf{1.} $v = \sqrt{2eU/m} = \qty{8.4e7}{m/s}$; $p = \qty{7.6e-23}{kg.m/s}$; $\lambda =
\qty{8.7}{pm}$।

\textbf{2.} कुछ प्रतिशत की त्रुटियाँ: आकलनों के लिए स्वीकार्य ($\lambda =
\qty{8.6}{pm}$ ठीक-ठीक)।

\textbf{3.} \qty{4.8}{ns}।

\textbf{4.} $\lambda/a = \qty{3e-8}{rad}$: \qty{12}{nm} --- किसी पिक्सेल से $10^4$ गुना
छोटा।

\textbf{5.} $\Delta\theta = \Delta p_y/p \sim \hbar/2ap = \lambda/4\pi a$: वही कोटि।

\textbf{6.} $\tau = \qty{1.6}{ms} \gg \qty{4.8}{ns}$: अपरिवर्तित।

\textbf{7.} किसी एकसमान क्षेत्र के लिए ठीक-ठीक $\dd\langle p_y\rangle/\dd t = eE$:
$\langle y\rangle$ चिरसम्मत परवलय मानता है; नली कोई
चिरसम्मत उपकरण है।

\textbf{8.} $v/2 = \qty{4.2e7}{m/s}$; केवल \omterm{prop:b2:dispersion-wave-packets:group}{समूह वेग} प्रेक्षणीय है;
किसी समतल तरंग की कला कोई संकेत नहीं ढोती (और ऊर्जा के मनमाने
शून्य पर निर्भर है)।

\textbf{9.} $\Delta p = \hbar/2\sigma_0$, $\Delta v = \hbar/2m\sigma_0$।

\textbf{10.} \qty{1.7e-16}{s}, \qty{1.7e-14}{s}, \qty{1.7e-8}{s}; प्रोटॉन \qty{3.2e-11}{s};
विषाणु \qty{20}{s}।

\textbf{11.} वह किसी \omterm{thm:b2:schrodinger-wave-functions:stationary}{स्थायी अवस्था} में है: विभव पुंज को थामे
रखता है; $|\varphi|^2$ हिलता नहीं।

\textbf{12.} $v = \qty{5.9e5}{m/s}$; दुगुना $t = \sqrt3\tau = \qty{2.9e-14}{s}$ पर:
\qty{17}{nm}, यानी कोई चौदह तरंगदैर्घ्य।

\textbf{13.} उनके \omterm{def:b2:schrodinger-wave-functions:psi}{तरंग फलन} समूचे क्रिस्टल पर फैले हैं: कोई
निश्चित संवेग, कोई निश्चित स्थिति नहीं।

\textbf{14.} विभेदन $\sim\lambda/\alpha = \qty{4}{pm}/0.01 = \qty{0.4}{nm}$: चुंबकीय
लेंसों के विपथन द्वारक को सीमित करते हैं, तरंगदैर्घ्य नहीं।

\textbf{15.} $\lambda$ की तुलना खेल में लगे आकारों से कीजिए (द्वारक, $V$
का मापक), और फैलाव काल की प्रेक्षण काल से।

\textbf{16.} \qty{0.38}{}, \qty{1.50}{}, \qty{3.38}{eV}; $\nu_{21} = \qty{2.7e14}{Hz}$
(\qty{1.1}{\micro m}), $\nu_{31} = \qty{7.3e14}{Hz}$ (\qty{410}{nm})।

\textbf{17.} कलादिक् का मापांक एक है; और $x = L/2$ पर।

\textbf{18.} $t = 0$ पर घनत्व बाएँ आधे में जमा होता है, और
$T_{21}/2$ पर दाएँ में।

\textbf{19.} $\langle x\rangle = L/2 - (16L/9\pi^2)\cos\omega_{21}t$: आयाम
$0.18L = \qty{0.18}{nm}$।

\textbf{20.} $\nu_{21}$ पर: वही वर्णक्रमी रेखा; किसी \omterm{thm:b2:schrodinger-wave-functions:stationary}{स्थायी अवस्था} का कोई
दोलायमान द्विध्रुव नहीं और वह स्थायी है; अध्यारोपण विकिरित करता हुआ
$\varphi_1$ तक उतर आता है।

\textbf{21.} निश्चित नहीं: $\langle E\rangle = (E_1 + E_2)/2 = \qty{0.94}{eV}$; कोई
मापन $E_1$ या $E_2$ देता है, हर एक आधी प्रायिकता के साथ।

\textbf{22.} यादृच्छिक: हर इलेक्ट्रॉन कहाँ गिरता है; यादृच्छिक नहीं: वह
प्रतिरूप $|\psi|^2$ जिसे गिरना बनाता है।

\textbf{23.} $j = |A|^2\hbar k/m$: प्रति इकाई लंबाई $|A|^2 = n$ इलेक्ट्रॉनों के साथ,
प्रति सेकंड $nv$ इलेक्ट्रॉन।

\textbf{24.} प्रथम कोटि ताकि $\psi(x,0)$ कोई पूर्ण वर्णन हो
जो भविष्य तय कर दे; सम्मिश्र ताकि किसी निश्चित ऊर्जा की
अकेली समय-निर्भरता $\eu^{-\iu Et/\hbar}$ हो और $E = p^2/2m$ निकल आए।

\textbf{25.} $\psi$ वह आयाम है जिसका वर्ग प्रायिकता है; वह
श्रोडिंगर का रैखिक समीकरण मानता है; \omterm{thm:b2:schrodinger-wave-functions:stationary}{स्थायी अवस्थाएँ} कोई कलादिक् और
नियत घनत्व ढोती हैं; कोई मुक्त पुंज $p/m$ पर चलता है और
$2m\sigma_0^2/\hbar$ में फैलता है; और एहरनफ़ेस्ट बड़ों के लिए न्यूटन लौटा देता है।
\end{solution}
